\documentclass{article}
\usepackage[T1]{fontenc}
\usepackage{lmodern}
\usepackage{graphicx}
\usepackage{amsmath,amssymb}
\usepackage[a4paper,margin=2cm]{geometry}
\usepackage[absolute,overlay]{textpos}
\setlength{\TPHorizModule}{1mm}
\setlength{\TPVertModule}{1mm}
\usepackage[indonesian]{babel}
\usepackage{hyperref}
\setlength{\emergencystretch}{2em}
% unit: Halaman 11

\begin{document}

% === Halaman 11 (llm) ===

\section*{1. Electronic Code Book (ECB)}

\begin{itemize}
    \item Setiap blok plainteks $P_i$ dienkripsi secara individual dan independen dari blok lainnya menjadi blok cipherteks $C_i$.
    \item Enkripsi: \[ C_i = E_K(P_i) \]
    Dekripsi: \[ P_i = D_K(C_i) \]
\end{itemize}

yang dalam hal ini, $P_i$ dan $C_i$ masing-masing blok plainteks dan cipherteks ke-$i$.

\end{document}
